
%BeginMC
\question

\begin{choices}
\correctchoice   %EndChoice
\choice  %EndChoice
\choice  %EndChoice
\choice  %EndChoice
\choice  %EndChoice
\end{choices}
%EndMC

%%%%%%%%%%%% Questions begin here. Copy and paste from template above.


%% Interspecific Competition, Niches, and Resource Partitioning %%

%BeginMC
\question
The ``role'' or ``job'' of a species in a community, represented by all of the resources needed by the species, is called 

\begin{choices}
	\correctchoice the niche.%EndChoice
	\choice the resource.%EndChoice
	\choice a predator.%EndChoice
	\choice a predator or prey.%EndChoice
	\choice a producer or consumer.%EndChoice
\end{choices}
%EndMC


%BeginMC
\question
Competition occurs between two different species in a community when

\begin{choices}
	\correctchoice both species need the same resources that is in limited supply.%EndChoice
	\choice both species are searching for mates.%EndChoice
	\choice one species is a predator of the other.%EndChoice
	\choice both species have partitioned resources.%EndChoice
	\choice one species is a parasite of the other.%EndChoice
\end{choices}
%EndMC

%BeginMC
\question
Resource partitioning

\begin{choices}
	\correctchoice increases the number of species that can coexist in a community because competition is reduced.%EndChoice
	\choice increases the number of individuals that can coexist in a population because competition is reduced.%EndChoice
	\choice reduces competition because the weaker competitor will be eliminated from the community.%EndChoice
	\choice reduces competition because the weaker competitor will be eliminated from the population.%EndChoice
	\choice reduces competition between individuals of the same species and between individuals of different species.%EndChoice
\end{choices}
%EndMC

%BeginMC
\question
Several species of warblers (a group of small colorful birds) in North America can live in the same tree only because they

\begin{choices}
	\choice created different types of nests within the tree.%EndChoice
	\choice eat different foods within the tree.%EndChoice
	\correctchoice partitioned resources to eliminate competition.%EndChoice
	\choice can fly to different heights within the tree.%EndChoice
	\choice out-competed other species of birds, such as sparrows.%EndChoice
\end{choices}
%EndMC

%BeginMC
\question
Several species of darters (a group of small fishes) live together on the gravel bottoms of small streams. They use similar resources in slightly different ways. The darters can live in the same habitat type  because they

\begin{choices}
	\choice have the same niches within the stream.%EndChoice
	\choice feed at different times of the day.%EndChoice
	\correctchoice partitioned the resources to reduce competition.%EndChoice
	\choice can swim at different heights above the gravel bottom.%EndChoice
	\choice out-competed other species of fishes, such as catfishes.%EndChoice
\end{choices}
%EndMC

%BeginMC
\question
Several species of anoles (a type of lizard) on Cuba can live in the same tree only because they

\begin{choices}
	\choice created new niches within the tree.%EndChoice
	\choice eat different foods within the tree.%EndChoice
	\correctchoice partitioned resources to eliminate competition.%EndChoice
	\choice can climb to different heights within the tree.%EndChoice
	\choice outcompeted other species of lizards, such as skinks.%EndChoice
\end{choices}
%EndMC

%BeginMC
\question
Species Old can nest anywhere in a certain kind of pine tree. Recently, due to global climate change, Species New has arrived from down south. Species New also nests in the same kind of pine tree but can only nest in the outermost parts of any branch (shown in darker color).  Species Old chases away Species New in the inner part of the tree but ignores Species New in the outer part of the tree. What will \emph{most likely} happen in the tree after a few years?

\begin{multicols}{2} 
	\begin{choices}
		\correctchoice Species New will nest in the outer part of the tree and Species Old will nest in the inner part of the tree.%EndChoice
		\choice Species Old will prevent Species New from making nests so Species New will die out in this location.%EndChoice
		\choice Species New will prevent Species Old from making nests so Species Old will die out in this location.%EndChoice
		\choice Both species will nest in the outer part of the tree in equal numbers but only Species Old will nest in the inner part of the tree.%EndChoice
		\choice Species Old will nest more often in the outer part of the tree.%EndChoice
		\choice Neither species will survive because they can't nest in the entire tree. %EndChoice
	\end{choices}
	\columnbreak
	{\centering \includegraphics[width=0.7\columnwidth]{exam4_partition_tree}\par }
\end{multicols}
%EndMC


%BeginMC
\question
Which of the following statements is consistent with the principle of competitive exclusion?  

\begin{choices}
	\correctchoice Even a slight reproductive advantage will eventually lead to the elimination of the less well adapted of two competing species. %EndChoice
	\choice Evolution tends to increase competition between related species.   %EndChoice
	\choice Bird species generally do not compete for nesting sites.  %EndChoice
	\choice The density of one competing species will have a positive impact on the population growth of the other competing species. %EndChoice
	\choice Two species with the same fundamental niche will exclude other competing species.   %EndChoice
\end{choices}
%EndMC

%BeginMC
\question
White-breasted nuthatches and brown creepers both eat insects that hide in the furrows of bark in hardwood trees. The brown creeper searches for insects by hunting from the bottom of the tree trunk toward the top, whereas the white-breasted nuthatch searches from the top of the trunk down. Their feeding areas rarely overlap on the tree trunk. These hunting behaviors best illustrate which of the following ecological concepts?

\begin{choices}
	\choice competitive exclusion%EndChoice
	\correctchoice resource partitioning%EndChoice
	\choice intraspecific competition%EndChoice
	\choice keystone species%EndChoice
	\choice interspecific competition%EndChoice
\end{choices}
%EndMC



%BeginMC
\question
Competitive exclusion means that

\begin{choices}
	\correctchoice one species will always out compete another species that occupies the same niche.  %EndChoice
	\choice two species will partition resources to eliminate competition. %EndChoice
	\choice one individual will always out compete another individual in a population. %EndChoice
	\choice one population will always exclude another another population of the same species in a community. %EndChoice
	\choice a species will have a smaller realized niche than a fundamental niche. %EndChoice
\end{choices}
%EndMC


%% Symbiosis %%%

%BeginMC
\question
A barnacle grows on a whale, doing it no harm. The movement of the whale exposes the barnacle to food. This is an example of

\begin{choices}
	\choice predation.%EndChoice
	\choice mutualism.%EndChoice
	\choice parasitism.%EndChoice
	\correctchoice commensalism.%EndChoice
	\choice resource partitioning.%EndChoice
\end{choices}
%EndMC

%BeginMC
\question
Which of the following statements about symbiotic relationships is true?

\begin{choices}
	\correctchoice A relationship that appears to be commensalism may in fact be mutualism or parasitism.%EndChoice
	\choice Both organisms are harmed in a parasitic relationship.%EndChoice
	\choice Both symbiotic organisms usually benefit from the relationship.%EndChoice
	\choice The most efficient type of parasite is one that kills its host.%EndChoice
	\choice Neither organism benefits in a symbiotic relationship.%EndChoice
\end{choices}
%EndMC

%BeginMC
\question
Malaria is a human disease that occurs when \textit{Plasmodium}, a single-celled organism, invades human blood cells.  The reproduction of \textit{Plasmodium} causes the red blood cells to burst, which decreases the ability of humans to carry oxygen. This is an example of 

\begin{choices}
	\choice commensalism.%EndChoice
	\choice mutualism.%EndChoice
	\correctchoice parasitism.%EndChoice
	\choice predation.%EndChoice
	\choice symbiosis.%EndChoice
\end{choices}
%EndMC


%BeginMC
\question
Cowbirds birds lay their eggs in the nests of other birds like cardinals. The cowbird chicks kick the cardinal chicks out of the nest, so the cardinal chicks die. The adult cardinals raise the cowbird chicks as their own. This is most likely an example of 

\begin{choices}
	\choice competition. %EndChoice
	\choice resource partitioning. %EndChoice
	\correctchoice parasitism. %EndChoice
	\choice symbiosis. %EndChoice
	\choice mutualism. %EndChoice
\end{choices}
%EndMC


%BeginMC
\question
Symbiosis is a longer-term ecological interaction

\begin{choices}
	\correctchoice between two or more species. %EndChoice
	\choice where one species benefits and another species is harmed %EndChoice
	\choice where both species benefit. %EndChoice
	\choice  where one species benefits and the other species neither benefits nor is harmed. %EndChoice
	\choice where two or more species partition resources. %EndChoice
\end{choices}
%EndMC

%BeginMC
\question
One form of elephantiasis occurs when \textit{Brugia malayyi}, a thread-like worm, invades human lymph nodes. The worm blocks the flow of lymph throughout the body, causing massive swelling of the lower body, especially legs and genitals. As described here, this is an example of

\begin{choices}
	\choice commensalism.%EndChoice
	\correctchoice parasitism.%EndChoice
	\choice mutualism.%EndChoice
	\choice predation.%EndChoice
	\choice symbiosis.%EndChoice
\end{choices}
%EndMC

%BeginMC
\question
Clownfish (like Nemo) live in sea anemones. The anemones will sting most organisms but not the clownfish. The clownfish keep debris out of the sea anemone and the waste from the clownfish provides food for the anemone. The stinging sea anemone provides a safe home for the clownfish. As described here, this is an example of

\begin{choices}
	\choice commensalism.%EndChoice
	\choice decomposition.%EndChoice
	\correctchoice mutualism.%EndChoice
	\choice parasitism.%EndChoice
	\choice symbiosis.%EndChoice
\end{choices}
%EndMC

%BeginMC
\question
Some fish, like the goby in the picture below, live inside sponges. Sponges are a type of living organism. The goby benefits from the shelter provided by the sponge. The sponge receives no apparent benefit nor is it harmed. As described here, this is an example of

\begin{multicols}{2} 
	\begin{choices}
		\correctchoice commensalism.%EndChoice
		\choice predation.%EndChoice
		\choice mutualism.%EndChoice
		\choice parasitism.%EndChoice
		\choice symbiosis.%EndChoice
	\end{choices}
	\columnbreak
	\hfill \includegraphics[width=4cm]{img/goby_sponge.jpg}
\end{multicols}
%EndMC



%% General Community %%

%BeginMC
\question
A community is correctly defined as having which of the following characteristics? \\
I.	Individuals inhabiting the same general area.\\
II.	Individuals belonging to the same species.\\
III. Individuals belonging to different species.\\
IV.	Individuals interacting with each other.

\begin{choices}
	\choice I and II only. %dontshuffle  %EndChoice
	\choice I and III only. %dontshuffle %EndChoice
	\choice I, II and IV only. %dontshuffle %EndChoice
	\correctchoice I, III and IV only. %dontshuffle %EndChoice
	\choice I and IV only. %dontshuffle %EndChoice
	\choice II and IV only. %dontshuffle %EndChoice
\end{choices}
%EndMC

%BeginMC
\question
Color that indicates to a predator that the prey is poisonous is known as

\begin{choices}
	\choice mimicry. %EndChoice
	\correctchoice aposomatic coloration. %EndChoice
	\choice cryptic coloration. %EndChoice
	\choice warning mimicry. %EndChoice
	\choice Batesian mimicry. %EndChoice
\end{choices}
%EndMC

%BeginMC
\question
Color that helps hide an animal in its environment is called

\begin{choices}
	\choice mimicry. %EndChoice
	\choice aposomatic coloration %EndChoice
	\correctchoice cryptic coloration. %EndChoice
	\choice warning coloration. %EndChoice
	\choice environmental coloration. %EndChoice
	\choice hideandseek coloration. %EndChoice
\end{choices}
%EndMC

%BeginMC
\question
A harmless fly looks like a dangerous wasp. The fly is using

\begin{choices}
	\choice camouflage. %EndChoice
	\correctchoice mimicry. %EndChoice
	\choice warning coloration. %EndChoice
	\choice cryptic coloration. %EndChoice
	\choice disruptive coloration. %EndChoice
\end{choices}
%EndMC


%BeginMC
\question
In a mature forest of oak, maple, and hickory trees, a disease causes a reduction in the number of acorns produced by oak trees. Which of the following would \emph{least} likely be a \emph{direct} result of this? 

\begin{choices}
	\correctchoice There might be a decrease in the number of maple seeds if the disease spreads to other trees in the forest. %EndChoice
	\choice There might be fewer squirrels because they feed on acorns produced by oak trees. %EndChoice
	\choice  There might be fewer mice and seed-eating birds because squirrels that normally eat acorns would eat more seeds and compete with the mice and birds.%EndChoice
	\choice There might be an increase in the number of hickory trees because the competition between hickory nuts and acorns for germination sites would be reduced or eliminated. %EndChoice
	\choice There might be fewer owls because they feed on baby squirrels, mice, and young seed-eating birds, whose populations would be reduced. %EndChoice
\end{choices}
%EndMC

%BeginMC
\question
Which block in the figure below shows the greatest species diversity? (Different shades, including white, represent different species. The small squares indicate relative abundance of individuals present.)

\begin{multicols}{3}
	\begin{choices}
		\correctchoice \includegraphics{diversity_mcE} %EndChoice
		\choice \includegraphics{diversity_mcA} %EndChoice
		\choice \includegraphics{diversity_mcB} %EndChoice
		\choice \includegraphics{diversity_mcC} %EndChoice
		\choice \includegraphics{diversity_mcD} %EndChoice
	\end{choices}
\end{multicols}
%EndMC

%BeginMC
\question
Which block in the figure below is most consistent with intermediate levels of disturbance in an ecosystem? (Different shades, including white, represent different species. The small squares indicate relative abundance of individuals present.)

\begin{multicols}{3}
	\begin{choices}
		\correctchoice \includegraphics{diversity_mcE} %EndChoice
		\choice \includegraphics{diversity_mcA} %EndChoice
		\choice \includegraphics{diversity_mcB} %EndChoice
		\choice \includegraphics{diversity_mcC} %EndChoice
		\choice \includegraphics{diversity_mcD} %EndChoice
	\end{choices}
\end{multicols}
%EndMC

%BeginMC
\question
In the experiment shown in the figure below, \textit{Balanus} was removed from the habitat shown on the left. Which of the following statements is a valid conclusion of this experiment?\\
\includegraphics[width=0.9\textwidth]{balanus.png}

\begin{choices}
	\choice \textit{Balanus} can survive only in the lower intertidal zone because it is unable to resist dehydration.%EndChoice
	\choice \textit{Balanus} is inferior to \textit{Chthamalus} in competing for space on rocks lower in the intertidal zone. %EndChoice
	\correctchoice The removal of \textit{Balanus} shows that the realized niche of \textit{Chthamalus} is smaller than its fundamental niche.%EndChoice
	\choice These two species of barnacle do not show competitive exclusion.%EndChoice
	\choice If \textit{Chthamalus} were removed, \textit{Balanus}'s fundamental niche would become larger.%EndChoice
\end{choices}
%EndMC


%BeginMC
\question
Naturally occurring fires cause occasional disturbance to an ecosystem. If humans stop these fires as soon as they are started, then

\begin{choices}
	\choice species diversity in the ecosystem will tend to increase.%EndChoice
	\correctchoice species diversity in the ecosystem will tend to decrease.%EndChoice
	\choice species diversity will tend to remain unchanged, but the types of species present may change.%EndChoice
	\choice the diversity of threatened and endangered species will be protected.%EndChoice
\end{choices}
%EndMC

%BeginMC
\question
Fire suppression by humans

\begin{choices}
	\choice will always result in an increase in species diversity in a given biome.%EndChoice
	\correctchoice can change the species composition within biological communities.%EndChoice
	\choice will result ultimately in sustainable production of increased amounts of forest products for human use.%EndChoice
	\choice is necessary for the protection of threatened and endangered forest species.%EndChoice
	\choice is a management goal of conservation biologists to maintain the healthy condition of forest communities.%EndChoice
\end{choices}
%EndMC

%BeginMC
\question
Forest mice eat seeds. Suppose a fire on the forest floor reduces the amount of seeds by 75\%. The mice survive the fire but what will happen to the size of the mouse population?

\begin{choices}
	\correctchoice It will decrease to about 75\% after a delay. %EndChoice
	\choice It will decrease rapidly by about 75\%. %EndChoice
	\choice It will decrease only a little as mice shift to another food source. %EndChoice
	\choice It will increase as the mice shift to another food source. %EndChoice
	\choice It will not be affected because mice survived the fire. %EndChoice
\end{choices}
%EndMC



%BeginMC
\question
Bobcats eat forest mice. The mice eat seeds. Suppose a fire on the forest floor reduces the amount of seeds by 75\%.  Before the fire, the population of mice sustained a bobcat population of 8. What will \emph{most likely} happen to the bobcat population as a result of the fire? The bobcat population will

\begin{choices}
	\correctchoice decrease to two after a delay. %EndChoice
	\choice decrease to two almost immediately. %EndChoice
	\choice decrease by two but most will survive. %EndChoice
	\choice increase as the bobcats shift to another food source. %EndChoice
	\choice not be affected because the mice survived the fire. %EndChoice
\end{choices}
%EndMC

%BeginMC
\question
Which of the following is \emph{not} directly concerned with community ecology?

\begin{choices}
	\choice herbivory%EndChoice
	\choice predation%EndChoice
	\correctchoice intraspecific competition%EndChoice
	\choice resource partitioning%EndChoice
	\choice symbiosis%EndChoice
\end{choices}
%EndMC

%BeginMC
\question
The difference between communities and ecosystems is that

\begin{choices}
	\choice communities include individuals of the same species; ecosystems include individuals of different species.%EndChoice
	\choice communities include individuals of the different species; ecosystems are individuals of the same species.%EndChoice
	\choice communities include interactions between individuals and their physical environment; ecosystems do \emph{not} include such interactions.%EndChoice
	\correctchoice communities do \emph{not} include interactions between individuals and the physical environment; ecosystems \emph{do} include such interactions.%EndChoice
	\choice communities are groups of individuals of the different species; ecosystems are groups of communities.%EndChoice
\end{choices}
%EndMC

%% - ORCA and OTTERS
%BeginMC
\question
After transient orca began to prey on sea otters, what happened to the kelp forest?

\begin{choices}
\choice The kelp abundance remained about the same. The kelp are not part of this community, so the kelp could not be affected.%EndChoice
\choice The kelp abundance remained about the same. Only the predator-prey interactions between Orca and otters were important, so nothing happened to the kelp.%EndChoice
\choice The kelp abundance increased. Otters dislodge kelp with their playful antics.  Because there were fewer otters, less kelp was dislodged so the kelp abundance increased.%EndChoice
\correctchoice The kelp abundance decreased. Sea otters eat urchins and urchins eat kelp.  Because there were fewer otters to eat urchins, the number of urchins increased.  The urchins therefore ate more kelp so the kelp abundance decreased.%EndChoice
\choice The kelp abundance decreased. Sea otters and kelp have a symbiotic relationship. Without the sea otters, the symbiotic relationship was broken and so the kelp abundance decreased.%EndChoice
\end{choices}
%EndMC


%BeginMC
\question
The population size of sea otters in southern Alaska decreased by an estimated 40,000 individuals between 1990--1996. The loss of sea otters reflects a greater change in the entire community of organisms in this area, due to human impacts. Which of the following statements is perhaps the best explanation for loss of sea otters?

\begin{choices}
\choice Global climate change caused the fish-eating pods of Orca to leave the area. They were replaced by the transient Orca that eat mammals.%EndChoice
\correctchoice Humans overfished the area, causing a decline in seal abundance. The transient Orca had to switch from seals to otters.%EndChoice
\choice Humans overfished other areas, causing the mammal-eating transient Orca to move to the area where the otters lived.%EndChoice
\choice Industrial activity is causing global climate change. The changing climate is affecting the interactions between all predators and their prey, including Orca and otters.%EndChoice
\choice Too much pollution, including oil, is entering the sea. This is probably creating a dead zone in the area. Once the other animals like fishes and urchins leave the dead zone, the otters starve to death.%EndChoice
\end{choices}
%EndMC

%BeginMC
\question
If humans stopped fishing along the coast of Alaska, what will probably happen to the kelp forest and why?

\begin{choices}
\correctchoice The kelp abundance will increase because the sea urchin abundance will probably decrease.%EndChoice
\choice The kelp abundance will remain about the same. The kelp are not part of this community, so the kelp would not be affected.%EndChoice
\choice The kelp abundance will decrease because more fish will be present to eat the kelp.%EndChoice
\choice The kelp abundance will decrease. Sea otters eat urchins and urchins eat kelp. Because there were fewer otters to eat urchins, the number of urchins increased. The urchins therefore ate more kelp so the kelp abundance decreased.%EndChoice
\choice The kelp abundance will increase because of the symbiotic relationship between kelp and fishes. Fishes are necessary for kelp to reproduce.%EndChoice
\end{choices}
%EndMC



%BeginMC
\question
The dispersion pattern for individuals in across the entire metapopulation shown in question \ref{ques:metapopulation} is most likely

\begin{choices}
	\correctchoice clumped. %EndChoice
	\choice uniform. %EndChoice
	\choice random. %EndChoice
	\choice even. %EndChoice
	\choice gene flow. %EndChoice
\end{choices}
%EndMC


%BeginMC
\question
While observing a population of lizards during one week, you notice that the number of adults has increased and is higher at the end of the week than at the start. The number of juveniles is unchanged. The most likely explanation for such an observation would be 

\begin{choices}
\correctchoice increased immigration.  %EndChoice
\choice decreased death rate. %EndChoice
\choice increased emigration.  %EndChoice
\choice decreased immigration. %EndChoice
\choice decreased emigration.  %EndChoice
\choice increased birth rate.  %EndChoice
\end{choices}
%EndMC
	

%BeginMC
\question
The attribute that most influence population density and rates of immigration and emigration among habitat patches in a metapopulation are

\begin{choices}
\correctchoice patch size, patch isolation and habitat quality. %EndChoice
\choice patch size, patch isolation and species diversity. %EndChoice
\choice patch isolation, habitat quality, and size of source patches. %EndChoice
\choice patch size, habitat quality, and isolation of sink patches. %EndChoice
\choice habitat quality, size of sink patches and isolation of source patches. %EndChoice
\end{choices}
%EndMC

%BeginMC
\question
The basic model of metapopulations views the number of suitable habitat patches occupied at a given point in time as a balance between

\begin{choices}
\correctchoice recolonization and local extinction. %EndChoice
\choice dispersal and birth rate. %EndChoice
\choice birth rate and local extinction. %EndChoice
\choice recolonization and degree of isolation. %EndChoice
\choice resource quality and degree of isolation. %EndChoice
\end{choices}
%EndMC



%BeginMC
\question
A small population of white-footed mice has the same intrinsic rate of increase (\textit{r}) as a large population.  If everything else is equal, 

\begin{choices}
\correctchoice the large population will add more individuals per unit time.  %EndChoice
\choice the small population will add more individuals per unit time. %EndChoice
\choice the two populations will add equal numbers of individuals per unit time.%EndChoice
\choice the growth trajectories of the two populations will proceed in opposite directions. %EndChoice
\choice the number of individuals added will depend on which population has the higher birth rate. %EndChoice
\end{choices}
%EndMC

%% Possible choice to include in above question.
\choice the number of individuals added will depend on the value of K. %EndChoice


%BeginMC
\question
Which of the following areas of study focuses on the exchange of energy and nutrients among individuals of different species?

\begin{choices}
\correctchoice Ecosystem ecology. %EndChoice
\choice Population Ecology %EndChoice
\choice Evolutionary Ecology %EndChoice
\choice Community Ecology %EndChoice
\choice Behavioral Ecology %EndChoice
\end{choices}
%EndMC


%BeginMC
\question
A patch in a population has the following four values. 

\qquad\begin{tabular}{@{}ll@{}}
$b = 0.04$ & $i = 0.06$\\
$d = 0.08$ & $e = 0.04$\\
\end{tabular}

Based only on these values (ignoring habitat quality, patch size, etc.), the patch is most likely to be a 

\begin{choices}
\correctchoice  sink.%EndChoice
\choice  source. %EndChoice
\end{choices}
%EndMC


%BeginMC
\question
Competition for limited resources may influence gene flow by causing

\begin{choices}
\correctchoice some individuals to migrate to another population.%EndChoice
\choice  some individuals to partition resources, reducing intraspecific competition.%EndChoice
\choice  an increased death rate, thereby reducing gene flow.%EndChoice
\choice  some individuals to move to a different source patch in the same population.%EndChoice
\choice  some individuals to evolve a new trait to be better competitors.%EndChoice
\end{choices}
%EndMC


%BeginMC
\question
Given the following values and assuming exponential population growth, what will be the total population size after one generation has elapsed? Round to nearest whole individual.\vspace{1ex}

\begin{minipage}{1em}
\end{minipage}%
\hspace{2em}%
\begin{minipage}{1.5in}
$b = 0.09$ \\
$d = 0.12$ \\
$N = 18,657$\\
$K = 15,000$\\
\end{minipage}

\begin{choices}
\choice $15,000$%EndChoice
\choice $560$%EndChoice
\correctchoice $18,097$%EndChoice
\choice $19,217$%EndChoice
\choice $3,657$%EndChoice
\end{choices}
%EndMC


%BeginMC
\question
Which of the following statements about ecology is \emph{incorrect}? 

\begin{choices}
\choice Ecologists may study populations and communities of organisms. %EndChoice
\correctchoice Ecology is a discipline that is independent from natural selection and evolutionary history. %EndChoice
\choice Ecology spans increasingly comprehensive levels of organization, from individuals to ecosystems. %EndChoice
\choice Ecology includes the study of the interactions between biotic and abiotic aspects of the environment. %EndChoice
\choice Ecologists may study the behavioral interactions between two species.%EndChoice
\end{choices}
%EndMC

%BeginMC
\question
Which of the following is most likely to contribute to density-dependent regulation of populations?

\begin{choices}
\choice the removal of toxic waste by decomposers%EndChoice
\correctchoice intraspecific competition for nutrients%EndChoice
\choice earthquakes%EndChoice
\choice floods%EndChoice
\choice fires%EndChoice
\end{choices}
%EndMC

%BeginMC
\question
Often the growth cycle of one population has an effect on the cycle of another. As moose populations increase, for example, wolf populations also increase. If you consider the logistic equation for the wolf population, which of the factors accounts for the effect of the moose population on the size of the wolf population?

\begin{choices}
\choice $r$%EndChoice
\choice $N$%EndChoice
\choice $rN$%EndChoice
\correctchoice $K$%EndChoice
\choice $\Delta t$%EndChoice
\choice $\Delta N$%EndChoice
\end{choices}
%EndMC

%BeginMC
\question
In the experiment shown in the figure below, \textit{Balanus} was removed from the habitat shown on the left. Which of the following statements is a valid conclusion of this experiment?\\
\includegraphics[width=0.9\textwidth]{exam3_balanus.png}

\begin{choices}
	\choice \textit{Balanus} can survive only in the lower intertidal zone because it is unable to resist dehydration.%EndChoice
	\choice Balanus is inferior to \textit{Chthamalus} in competing for space on rocks lower in the intertidal zone. %EndChoice
	\correctchoice The removal of \textit{Balanus} shows that the realized niche of \textit{Chthamalus} is smaller than its fundamental niche.%EndChoice
	\choice These two species of barnacle do not show competitive exclusion.%EndChoice
	\choice If \textit{Chthamalus} were removed, \textit{Balanus}'s fundamental niche would become larger.%EndChoice
\end{choices}
%EndMC


